%!TEX root = ./../../springer-mono-format/book.tex
%!TEX root = ./../../thesis-trier.tex

In the present section, we investigate automata all of whose strongly connected components are blocks. These automata arose from the study of solvable automata, which we introduce later in \cref{prelim:subclasses:solvable}. We observed that several properties of solvable automata extend to this more general class.

First, \cref{prelim:thm:SCCblocks-rank_r} gives a bound on the length of the shortest words of rank~$r$. \Cref{prelim:classes:block:prop:sync-into-bound} gives bounds on the length of the shortest words mapping the entire state set into an upper set. \Cref{prelim:classes:pto:prop:u-raised-to-n} gives words ending in certain maximal states. We derive several consequences in \cref{prelim:classes:block:cor:S-into-T,prelim:classes:block:cor:Q-into-S,prelim:classes:block:cor:Q-into-S-bound}. \Cref{prelim:classes:block:prop:char-sink-state} characterizes maximal states. Lastly, in \cref{prelim:class:block:cor:sync_trap_state}, we show that every synchronizing state is a trap state.


\begin{theorem} \label{prelim:thm:SCCblocks-rank_r}
	Let $\semiautomat{A}{Q}{\Sigma}{\delta}$ be an $n$-state automaton all of whose strongly connected components are blocks that has a word of rank~$r$.
	Then there exists a word of length
	at most $n - r$ whose rank is at most $r$,
	and this bound is tight.
	
	Furthermore, let $u \in \bigcap_{a \in \Sigma}\Sigma^* a \Sigma^*$.
	Then $u^{n - r}$
	has rank at most $r$, and $\delta(Q, u^{n - 1})$
	is the union of all sink components of~$\mathcal A$.
\end{theorem}
\begin{proof}
	Suppose $\mathcal A$ admits a word of rank $r$.
	We can assume $n > r$.
	By Proposition~\ref{prelim:cor:composition-series-maximal-states}
	there exists a sequence of upper sets 
	\begin{equation}\label{eqn:sequence}
		S_1 \subsetneq \cdots \subsetneq S_m = Q
	\end{equation}
	where~$S_1$ equals the union of all sink components
	and $S_{i} \setminus S_{i-1}$ for $i \in \{2,\ldots,m\}$ 
	are all the strongly connected components that are not sink components.
	Let $i \in \{2,\ldots,m\}$
	and set $S = S_{i} \setminus S_{i-1}$.
	As $S$ is not a sink component,
	we have $\delta(S, x_i) \setminus S \ne \emptyset$
	for some $x_i \in \Sigma$,
	which gives $\delta(S, x_i) \cap S = \emptyset$
	as $S$ is a block.
	As $S_i$ is an upper set, 
	we have $\delta(S, x_i) \subseteq S_i = S_{i-1} \cup S$ 
	and so $\delta(S, x_i) \subseteq S_{i-1}$.
	Hence, we find
	\[
	\delta(S_i, x_i) = \delta(S_{i-1} \cup S, x_i) = \delta(S_{i-1}, x_i) \cup \delta(S, x_i) \subseteq S_{i-1}.
	\]
	This also gives $\delta(S_i, u) \subseteq S_{i-1}$ for all $i \in \{2,\ldots,m\}$
	since $\delta(S_i, a) \subseteq S_i$ and $\delta(S_{i-1}, a) \subseteq S_{i-1}$
	for all $a \in \Sigma$.
	Now, inductively,  
	from $\delta(S_i, x_i) \subseteq S_{i-1}$
	and  $\delta(S_i, u) \subseteq S_{i-1}$ for all $i \in \{2,\ldots,m\}$,
	we obtain that
	for all $i \in \{2,\ldots,m\}$ we have
	\begin{equation*}\label{eqn:word-each-letter}
		\delta(Q, x_m \cdots x_i) \subseteq S_{i-1}
		\quad\text{ and }\quad
		\delta(Q, u^{m - (i - 1)}) \subseteq S_{i-1}.
	\end{equation*}
	Since the sequence from Equation~\eqref{eqn:sequence}
	strictly decreases, we get $|S_i| \le n - (m - i)$
	for $i \in \{1,\ldots,m\}$.
	%  Now, for a sink component $S \subseteq Q$ and $x \in \Sigma$
	%  we have $\delta(S, x) \subseteq S$
	%  and as it is a block this gives $\delta(S, x) = S$, i.e., $x$
	%  permutes $S$. 
	As~$S_1$ is a union of sink components that are blocks, 
	from Equation~\eqref{eqn:closed-block} we get $S_1 = \delta(S_1, u) \subseteq \delta(Q, u)$
	for each $u \in \Sigma^*$, which gives $S_1 \subseteq \delta(Q,u^{n-1})$  and $|S_1| \le r$.
	As $|S_1| \le r$ there exists a maximal $j \in \{1,\ldots,m\}$
	such that $|S_j| \le r$. 
	We have $j \ne m$
	and $r < |S_{j+1}| \le n - (m - (j+1))$.
	So we get $m - j - 1 < n - r$.
	Now, as $\delta(Q, x_m \cdots x_{j+1}) \subseteq S_j$
	and $\delta(Q, u^{m - j}) \subseteq S_{j}$
	the words $x_m \cdots x_{j+1}$ 
	and $u^{m - j}$
	of length $m - j \le n - r$ and $|\Sigma|(m-j)$, respectively,
	have rank at most~$r$,
	and the same follows for $u^{n- r}$.
	Lastly, $\delta(Q, u^{m - 1}) \subseteq S_1$
	gives $\delta(Q, u^{n - 1}) \subseteq S_1$ as $m - 1 \le n - 1$
	and so $\delta(Q, u^{n - 1}) = S_1$.
	
	For tightness, consider the automata from the proof of \cref{prelim:prop:unary_bound}.
\end{proof}

\begin{remark}
	\Cref{prelim:automata:cor:composition-series-maximal-states} 
	can be used to give an inductive
	proof of \cref{prelim:thm:SCCblocks-rank_r}. 
	We can assume~$\mathcal A$ admits a word $u \in \Sigma^*$ of rank~$r$
	with $r < n$ for $n = |Q|$.
	If~$Q$ equals the set of maximal states,
	then~$Q$ equals a union of sink components.
	By assumption, these sink components are blocks.
	By \cref{eqn:closed-block} 
	each letter permutes each sink component.
	Thus, $\mathcal A$ is a permutation automaton.
	Now, for  $r < |Q|$, this cannot be the case.
	Then \cref{prelim:automata:cor:composition-series-maximal-states}
	yields existence of a proper and non-empty
	upper set $S \subseteq Q$
	such that $\mathcal A / {\sim_S}$
	is strongly trap-connected.
	Since $|\delta(S, u)| \le r$,
	inductively, there exists $v \in \Sigma^*$
	with $|v| \le \max\{ |S| - r, 0 \}$
	such that $|\delta(S, v)| \le r$.
	Furthermore, we have $\delta(Q \setminus S, x) \cap S \ne \emptyset$
	for some~$x \in \Sigma$ since~$Q \setminus S$ is not a sink component
	by \cref{prelim:fact:factor-strongly-trap}. 
	Thus, as $Q \setminus S$ is a block,
	we find $\delta(Q \setminus S, x) \subseteq S$.
	Hence $\delta(Q, x) \subseteq S$.
	Then $|\delta(Q, xv)| \le r$
	and $|xv| \le 1 + \max\{ |S| - r, 0 \} \le n - r$.
	For $w \in \bigcap_{a \in \Sigma} \Sigma^* a \Sigma^*$, we can also inductively
	suppose $|\delta(S, w^{\max\{|S|-r,0\}})| \le r$
	and, similarly as above, conclude
	$|\delta(Q, w^{n - r})| \le r$.
\end{remark}

\begin{proposition} \label{prelim:classes:block:prop:sync-into-bound}
	Let $\semiautomat{A}{Q}{\Sigma}{\delta}$ be an automaton all of whose strongly connected components are blocks and $S \subseteq Q$ be an upper set.
	Suppose $\delta(Q, u) \subseteq S$ for some $u \in \Sigma^*$.
		Then there exists $v \in \Sigma^*$
		with $\delta(Q, v) \subseteq S$ and $|v| \le |Q| - |S|$,
		and this bound is tight.
\end{proposition}
\begin{proof}
	We may assume $Q \ne \emptyset$; then $\emptyset \ne \delta(Q,u) \subseteq S$, so $S \ne \emptyset$.
	We induct on~$|Q|$.
	
	If $S = Q$, then $v = \varepsilon$ works. So assume $S \subsetneq Q$.
	Then $S$ is a non-empty proper upper set, so we may choose an inclusion-maximal
	proper non-empty upper set $T \subseteq Q$ and put $C := Q \setminus T$.
	By \cref{prelim:automata:prop:TsetminusS-strong-conn,prelim:automata:cor:SsetminusT},
	$C$ is a strongly connected component, hence a block by assumption.
	
	As $T$ is an upper set, $\mathcal A_{|T}$ is well defined and its strongly connected
	components are strongly connected components of~$\mathcal A$, hence blocks.
	Moreover $T \cap S$ is an upper set of $\mathcal A_{|T}$ and
	$\delta(T,u) \subseteq T \cap \delta(Q,u) \subseteq T \cap S$.
	Since $|T| < |Q|$, induction gives $w \in \Sigma^*$ with
	$\delta(T,w) \subseteq T \cap S$ and $|w| \le |T| - |T \cap S| = |T \setminus S|$.
	
	\emph{Case 1: $C \subseteq S$.} Then $Q \setminus S \subseteq T$, so
	$|w| \le |T \setminus S| = |Q \setminus S| = |Q| - |S|$.
	As $S$ is an upper set, $\delta(C,w) \subseteq S$, whence
	$\delta(Q,w) = \delta(C,w) \cup \delta(T,w) \subseteq S$. Take $v = w$.
	
	\emph{Case 2: $C \not\subseteq S$.} 	
	First, $\delta(C,a) \cap T \ne \emptyset$ for some
	$a \in \Sigma$: otherwise $\delta(C,a) \subseteq C$ for all $a \in \Sigma$, so $C$ and
	hence $C \cap S$ are upper sets; the latter is non-empty because
	$\emptyset \ne \delta(C,u) \subseteq C \cap S$, and it is a proper subset of~$C$,
	so $T \subsetneq T \cup (C \cap S) \subsetneq Q$ contradicts the maximality of~$T$.
	Fix such an~$a$. Since $C$ is a block, $\delta(C,a) \cap C = \emptyset$, i.e.
	$\delta(C,a) \subseteq T$, and $\delta(T,a) \subseteq T$; thus $\delta(Q,a) \subseteq T$
	and $\delta(Q,aw) \subseteq \delta(T,w) \subseteq S$.
	Finally $Q \setminus (S \cup T) = C \setminus S \ne \emptyset$, so
	\[
	|Q| - |S| = |Q \setminus (S \cup T)| + |T \setminus S|
	\ge 1 + |T \setminus S| \ge 1 + |w| = |aw| .
	\]
	Take $v = aw$.
	
	Tightness follows from the automata in \cref{prelim:fig:unary_tight}
	with $S = \{q_{n-r}, \ldots, q_n\}$.
\end{proof}

\Cref{prelim:classes:block:prop:sync-into-bound} 
might fail if~$S$ is not an upper set. For example, 
in the automaton from \cref{prelim:classes:rem:unary-seq-aut},
the shortest word mapping all states into $\{1, n\}$
is $a^{n - 1}$.





\begin{proposition} \label{prelim:classes:pto:prop:u-raised-to-n}
	Let $\semiautomat{A}{Q}{\Sigma}{\delta}$ be an automaton all of whose strongly connected components are blocks and $S \subseteq Q$. Then, for $u \in \bigcap_{a \in \Sigma} \Sigma^* a \Sigma^*$
	and $n \ge |Q| - 1$	we have
	\[
	\delta(S, u^{n}) \subseteq \{ q \in \upperset{S} \mid q \text{ is a maximal state } \}. 
	\]
\end{proposition}
\begin{proof}
	Let $q \in S$. Clearly, $\delta(q, u^n) \in \upperset{S}$.
	By finiteness,
	we have $\delta(q, u^{i + p}) = \delta(q, u^i)$
	for some $0 \le i < |Q|$ and $p > 0$.
	Let $T$ be the strongly connected component~$T$
	with $\delta(q, u^i) \in T$.
	We have 
	\[
	\delta(q, u^{i+p}) \eqreach \delta(q, u^{i+p-1}) \eqreach \cdots \eqreach \delta(q, u^{i})
	\]
	Thus $\delta(T, u) \cap T \ne \emptyset$.
	Since~$T$ is a block, we find $\delta(T, u) = T$.
	Because~$u$ contains each letter,
	this yields $\delta(T, a) = T$
	for each $a \in \Sigma$
	and so~$T$ is a sink component.
	Hence, $\delta(q, u^i)$ is maximal,
	and so $\delta(q, u^n)$ is maximal.	
\end{proof}

The containment in \cref{prelim:classes:pto:prop:u-raised-to-n} 
can be proper. For example, this is the case for $S = \{q_1\}$ and $u = ab$ in 
the automaton whose strongly connected components are blocks from \cref{prelim:fig:rho-aut-left}

\begin{corollary}  \label{prelim:classes:block:cor:S-into-T}
	\label{prelim:classes:block:cor:equalR}
	Suppose $\semiautomat{A}{Q}{\Sigma}{\delta}$ is an automaton all of whose strongly connected components are blocks.
	Let $S \subseteq Q$ and $R = \{ q \in \upperset{S} \mid \text{$q$ is a maximal state} \}$.
	Suppose $R \subseteq S$.
	
	\begin{enumerate}
	\item \label{prelim:classes:block:cor:S-into-T:2} For $T \subseteq Q$ we have $\delta(S, u) \subseteq T$ for some $u \in \Sigma^*$ if and only if $R \subseteq T$.
		
	\item \label{prelim:classes:block:cor:S-into-T:1} For $u \in \bigcap_{a \in \Sigma} \Sigma^* a \Sigma^*$
	and $n \ge |Q| - 1$	we have
	$
	\delta(S, u^{n}) = R.
	$
	\end{enumerate}	
\end{corollary}
\begin{proof}
 \ref{prelim:classes:block:cor:S-into-T:2}: First, let $\delta(S, u) \subseteq T$
 for some $u \in \Sigma^*$.
 For $R \subseteq S$, the set $R$ is the union of the sink components contained in~$S$.
 By \cref{eqn:closed-block}, $\delta(S, u)$
 also contains these sink components.
 Thus $R$ is contained in $\delta(S, u)$ and therefore in~$T$.
 
 Conversely, if $R \subseteq T$, 
 by \cref{prelim:classes:pto:prop:u-raised-to-n} 
 there exists a word mapping~$Q$ into $R$,
 and thus into~$T$.
 
 \ref{prelim:classes:block:cor:S-into-T:1}: Set $T = \delta(S, u^n)$.
 We have $\delta(S, u^n) \subseteq T$.
 Thus, by \cref{prelim:classes:block:cor:S-into-T:2}, we get $R \subseteq T$.
 By \cref{prelim:classes:pto:prop:u-raised-to-n} we have $T \subseteq R$.
\end{proof}

Note that $R \subseteq S$ holds in \cref{prelim:classes:block:cor:S-into-T} if~$S$ is an upper set.

\Crefitem{prelim:classes:block:cor:S-into-T}{prelim:classes:block:cor:S-into-T:2} 
generalizes the 
statement from \cref{prelim:thm:SCCblocks-rank_r} that words which contain each letter applied $|Q| - 1$ times map the whole state set onto the set of maximal states.

\begin{corollary} \label{prelim:classes:block:cor:Q-into-S}
	Let $\semiautomat{A}{Q}{\Sigma}{\delta}$ be an automaton all of whose strongly connected components are blocks.
	For $S \subseteq Q$ there exists $u \in \Sigma^*$ 
	such that $\delta(Q, u) \subseteq S$
	if and only if $T \subseteq S$
	for each sink component $T \subseteq Q$ of~$\mathcal A$.	
\end{corollary}
\begin{proof}
	This follows from \cref{prelim:classes:block:cor:S-into-T}
	applied to~$Q$ and~$S$ (since in this case, $R$ equals the union of all sink components).
\end{proof}

\begin{corollary} \label{prelim:classes:block:cor:Q-into-S-bound}
	Let $\semiautomat{A}{Q}{\Sigma}{\delta}$ be an automaton all of whose strongly connected components are blocks and $S \subseteq Q$.
Suppose $\delta(Q, u) \subseteq S$ for some $u \in \Sigma^*$.
Then there exists $v \in \Sigma^*$
with $\delta(Q, v) \subseteq S$ and $|v| \le |Q| - 1$.		
\end{corollary}
\begin{proof}
	Let~$R$ be the union of all sink components in~$\mathcal A$.
	By \cref{prelim:classes:pto:prop:u-raised-to-n} 
	there exists a word mapping~$Q$ into~$R$.
	Thus, by \cref{prelim:classes:block:prop:sync-into-bound},
	we have $\delta(Q, u) \subseteq R$
	for some $u \in \Sigma^*$ with $|u| \le |Q| - |R| \le |Q| - 1$.
	By \cref{prelim:classes:block:cor:Q-into-S},
	we have $R \subseteq S$
	and so~$u$ maps~$Q$ into~$S$.	
\end{proof}

Next, we characterize the maximal states.

\begin{proposition} \label{prelim:classes:block:prop:char-sink-state}
	Let $\semiautomat{A}{Q}{\Sigma}{\delta}$ be an automaton all of whose strongly connected components are blocks.
	A state $q \in Q$ is contained in a sink component
	if and only if for all $x \in \Sigma$ we have $\delta(q, x^i) = q$
	for some $i \in \{1,\ldots,|Q|\}$.
\end{proposition}
\begin{proof}
	Let $S \subseteq Q$ be the strongly connected component 
	with $q \in S$.
	First, assume~$S$ is a sink component.
	By Equation~\eqref{eqn:closed-block}, every letter $x \in \Sigma$
	permutes $S$ and so $\delta(q, x^i) = q$
	for some $i \in \{1,\ldots,|S|\}$.
	Next, assume~$S$ is not a sink component.
	Then, as~$S$ is a block,
	we have $\delta(S, x) \cap S = \emptyset$
	for some $x \in \Sigma$.
	As~$S$ is a strongly connected component,
	$\delta(q, x) \notin S$ yields $\delta(q, x^i) \notin S$,
	and so $\delta(q, x^i) \ne q$,
	for all $i \ge 1$.
\end{proof}

\begin{proposition} \label{prelim:class:block:cor:sync_trap_state} \label{constr_sync:prop:sink-blocks-sync-trap}
	Let $\semiautomat{A}{Q}{\Sigma}{\delta}$ be an automaton all of whose sink components are blocks.
	Then each synchronizing state is a trap state.
\end{proposition}
\begin{proof}
	By \cref{prelim:thm:SCC_sync_words}, the set of synchronizing states~$S$
	is a sink component.
	Thus $\mathcal A_{|S}$ is a synchronizing permutation automaton. 
	This implies $|S| = 1$.
\end{proof}

So, by \cref{prelim:class:block:cor:sync_trap_state} 
an automaton in which all sink components are blocks is synchronizable
if and only if it has a unique trap state that is reachable from every other state. 
